% ===========================================================================
% Poste 10 — Allocation-logement
% ===========================================================================
\poste{10}{Programme Allocation-logement}{QC\_al}

\pdesc Aide mensuelle de la Société d'habitation du Québec, administrée par
Revenu Québec, pour les ménages à faible revenu qui consacrent une part
excessive de leur revenu au logement. Sont admissibles les ménages dont un
adulte a 50~ans ou plus, ou qui comptent au moins un enfant à charge.

\pobj Réduire le fardeau du logement des ménages à faible revenu --- aînés
et familles --- sans passer par le logement social.

\pregle Le programme réel calcule un \emph{taux d'effort} : la part du
revenu consacrée au loyer. Le calculateur ne demandant pas le loyer, le
modèle l'\emph{impute} selon la composition du ménage (voir Notes). Le
calcul, fidèle au calculateur officiel :
\begin{enumerate}
  \item admissibilité : un adulte de 50~ans ou plus, \emph{ou} $n \ge 1$ ;
  \item taux d'effort, arrondi à quatre décimales :
        $\varepsilon = (loyer \times 12) / \RAL$, où $\RAL$ est le revenu
        ajusté de l'équation~\ref{eq:ral} ;
  \item montant mensuel selon le palier d'effort :
        \[
        m = \begin{cases}
        0 & \varepsilon < \pc{30} \\
        \D{100} & \pc{30} \le \varepsilon < \pc{50} \\
        \D{150} & \pc{50} \le \varepsilon < \pc{80} \\
        \D{170} & \varepsilon \ge \pc{80}
        \end{cases}
        \]
  \item réduction \emph{dollar pour dollar} du revenu excédant le seuil
        $s$ (selon la composition) :
        \[
        AL = \pos{12 \times m - \pos{\RAL - s}}
        \]
\end{enumerate}

\pparams

\begin{center}
\small
\begin{tabular}{l r r}
\toprule
Paramètre & 2025 & 2026 \\
\midrule
Seuil --- 1 adulte, 0 enfant ($s$)        & \D{22400} & \D{22900} \\
Seuil --- 2 adultes, 0 enfant             & \D{31500} & \D{32100} \\
Seuil --- 1 ad.\ 1--2 enf.\ ; 2 ad.\ 1 enf. & \D{38700} & \D{39500} \\
Seuil --- 1 ad.\ 3+ enf.\ ; 2 ad.\ 2+ enf.  & \D{44600} & \D{45500} \\
Montant mensuel --- effort 30 à 50 \%     & \D{100}   & \D{100} \\
Montant mensuel --- effort 50 à 80 \%     & \D{150}   & \D{150} \\
Montant mensuel --- effort 80 \% et +     & \D{170}   & \D{170} \\
Âge d'admissibilité (sans enfant)         & 50 ans    & 50 ans \\
\bottomrule
\end{tabular}
\end{center}

Loyers mensuels imputés par le modèle (non indexés en 2026) : 1 adulte sans
enfant \D{856} ; 1 adulte avec 1 enfant ou couple sans enfant \D{1002} ;
1 adulte 2 enfants ou couple 1 enfant \D{1131} ; compositions plus grandes
\D{1380}.

\pnotes (1)~Le loyer imputé est un choix de modélisation hérité du
calculateur officiel, non une règle du programme --- le programme réel
utilise le loyer déclaré. (2)~Pour les aînés, le revenu $\RAL$ retranche du
revenu familial net une fraction des pensions (facteur $f$ de
l'équation~\ref{eq:ral}, extrait du calculateur officiel). (3)~Le seuil
couple sans enfant 2026 (\D{32100}) diffère de \D{100} de celui publié par
la chaire en fiscalité (CFFP, \D{32200}) ; la valeur du calculateur
officiel est retenue.

% ============================================================================
% GÉNÉRÉ par scripts/exemples-guide.ts — NE PAS ÉDITER À LA MAIN.
% Régénérer : npm run exemples (chaque chiffre est vérifié contre le moteur).
% ============================================================================
\pexemples Le modèle \emph{impute} le loyer selon la composition du ménage, calcule
le taux d'effort (loyer annuel $\div$ revenu $\RAL$, arrondi à quatre
décimales), convertit ce taux en montant mensuel (0, \D{100},
\D{150} ou \D{170}), puis réduit le montant annuel
\emph{dollar pour dollar} au-delà du seuil de revenu.

\medskip\noindent\textbf{M6 --- famille monoparentale, 35~ans, revenu de travail de \D{35000}, un enfant : 3~ans (garde subventionnée, \D{2000}).}
Admissible (un enfant à charge). Loyer imputé (1~adulte, 1~enfant) :
\D{1002} par mois.
\begin{align*}
\text{effort} &= \frac{\num{1002} \times 12}{\num{33265}} = 0.3615 \ \Rightarrow\ \num{100}\,\$/\text{mois} \\
\text{réduction} &= \pos{\num{33265} - \num{38700}} = \num{0} \\
\text{allocation} &= \num{100} \times 12 - \num{0} = \num{1200}
\end{align*}
Allocation-logement : \D{1200} par année.

\medskip\noindent\textbf{M10 --- retraité vivant seul, 70~ans, revenu de pension privée de \D{20000}.}
Admissible (50~ans et plus). Le revenu $\RAL$ retranche du $\RFN$
une fraction des pensions (équation~\eqref{eq:ral}) : \num{29401.05}.
\begin{align*}
\text{effort} &= \frac{\num{856} \times 12}{\num{29401.05}} = 0.3494 \ \Rightarrow\ \num{100}\,\$/\text{mois} \\
\text{réduction} &= \pos{\num{29401.05} - \num{22400}} = \num{7001.05} \\
\text{allocation} &= \pos{\num{1200} - \num{7001.05}} = \num{0}
\end{align*}
La réduction dollar pour dollar gruge une partie du montant :
\D{0} par année.

\medskip\noindent\textbf{M13 --- couple de retraités, 68 et 66~ans, pensions privées de \D{12000} et \D{6000}.}
Couple d'aînés, loyer imputé de \D{1002} :
\begin{align*}
\text{effort} &= \frac{\num{1002} \times 12}{\num{40083.04}} = 0.3 \ \Rightarrow\ \num{100}\,\$/\text{mois} \\
\text{réduction} &= \pos{\num{40083.04} - \num{31500}} = \num{8583.04} \\
\text{allocation} &= \num{100} \times 12 - \num{8583.04} = \num{0}
\end{align*}
Allocation-logement : \D{0} par année.

\medskip\noindent\textbf{M7 --- couple, 45 et 44~ans, revenus de travail de \D{30000} et 0\,\$.}
Aucun adulte de 50~ans ou plus et aucun enfant : ménage
\emph{inadmissible}, allocation nulle — quel que soit son revenu.


\prefs Programme établi en vertu de la Loi sur la Société d'habitation du
Québec (RLRQ, c.~S-8), paramètres fixés par décret ; administré par Revenu
Québec ; Société d'habitation du Québec.
